% !TeX program = xelatex
\documentclass[a4paper,11pt]{article}
\usepackage[a4paper, top=2cm, bottom=2.5cm, left=2cm, right=2cm]{geometry}
\usepackage{fontspec}
\usepackage{polyglossia}
\setmainlanguage{portuguese}
\usepackage{microtype}
\usepackage{xcolor}
\usepackage{array}
\usepackage{tabularx}
\usepackage{booktabs}
\usepackage{enumitem}
\usepackage{fancyhdr}
\usepackage{tcolorbox}
\tcbuselibrary{skins,breakable}
\usepackage{hyperref}

\hypersetup{colorlinks=false,pdfborder={0 0 0},
  pdftitle={Analise SWOT -- Engenharia de Reservatorios 1}}

\definecolor{forcas}{HTML}{0F6E56}
\definecolor{forcasLight}{HTML}{E1F5EE}
\definecolor{opor}{HTML}{185FA5}
\definecolor{oporLight}{HTML}{E6F1FB}
\definecolor{fraq}{HTML}{854F0B}
\definecolor{fraqLight}{HTML}{FAEEDA}
\definecolor{amea}{HTML}{993C1D}
\definecolor{ameaLight}{HTML}{FAECE7}
\definecolor{axisColor}{HTML}{444441}
\definecolor{titleBg}{HTML}{2C2C2A}
\definecolor{pageGray}{HTML}{F1EFE8}

\pagestyle{fancy}
\fancyhf{}
\renewcommand{\headrulewidth}{0.4pt}
\fancyhead[L]{\small\textcolor{axisColor}{Análise SWOT}}
\fancyhead[C]{\small\textcolor{axisColor}{\textit{Engenharia de Reservatórios 1}}}
\fancyhead[R]{\small\textcolor{axisColor}{Rocélio Da Silva}}
\fancyfoot[C]{\small\textcolor{axisColor}{\thepage}}
\fancyfoot[R]{\small\textcolor{axisColor}{Perspectiva: Sucesso na aprovação}}

% Estilo das caixas com margens horizontais largas
\tcbset{
  swotbox/.style={
    enhanced, boxrule=0.6pt, arc=8pt,
    fonttitle=\bfseries\small, coltitle=white,
    top=6pt, bottom=6pt,
    left=18pt, right=18pt,            % margens horizontais internas generosas
    toptitle=4pt, bottomtitle=4pt,
    valign=top, height fill,          % preenche a altura, conteúdo alinhado ao topo
  }
}

% Lista com espaçamento elástico para preencher verticalmente a caixa
\newlist{swotlist}{itemize}{1}
\setlist[swotlist]{
  label=\textbullet,
  leftmargin=0.8em,
  labelsep=0.5em,
  itemsep=2pt plus 2pt,    % espaço entre itens pode aumentar para ocupar a altura
  parsep=0pt,
  topsep=2pt,
  partopsep=0pt,
  before={\setlength{\parindent}{0pt}\setlength{\parskip}{0pt}},
}

\raggedbottom

\begin{document}
\pagecolor{pageGray}

% ------------------ TÍTULO ------------------
\newsavebox{\titlebox}
\begin{lrbox}{\titlebox}
\begin{tcolorbox}[enhanced, colback=titleBg, colframe=titleBg, arc=10pt, boxrule=0pt,
  top=10pt, bottom=10pt, left=14pt, right=14pt]
  \centering
  {\Large\bfseries\color{white} ANÁLISE SWOT}\\[4pt]
  {\normalsize\color{white!80!titleBg}
    Unidade Curricular de \textbf{Engenharia de Reservatórios 1}}\\[2pt]
  {\small\color{white!60!titleBg}
    Docente: Dr. Geraldo Ramos $\bullet$ Discente: Rocélio Da Silva}
\end{tcolorbox}
\end{lrbox}

% Medir altura do título
\newlength{\titleheight}
\setlength{\titleheight}{\ht\titlebox}
\addtolength{\titleheight}{\dp\titlebox}
\addtolength{\titleheight}{4pt}

% Calcular altura restante para a grelha
\newlength{\gridheight}
\setlength{\gridheight}{\dimexpr\textheight - \titleheight - 2\baselineskip\relax}
\newlength{\rowheight}
\setlength{\rowheight}{\dimexpr0.5\gridheight - 4pt\relax}

% Imprimir título
\usebox{\titlebox}
\vfill

% ------------------ GRELHA 2x2 ------------------
\noindent
\begin{tabularx}{\textwidth}{@{} X @{\hspace{8pt}} X @{}}
  % Forças
  \begin{tcolorbox}[swotbox,
    title={FORÇAS (Strengths)},
    colback=forcasLight, colframe=forcas, colbacktitle=forcas,
    height=\rowheight]
    \begin{swotlist}
      \item \textbf{Professor Dr. Geraldo Ramos é muito experiente:} conhece bem a disciplina e já leccionou várias vezes.
      \item \textbf{Conteúdo didático disponibilizado regularmente:} todo o material é colocado à disposição no unimestre.
      \item \textbf{Aulas frequentes:} acontecem duas vezes por semana, ajudando a manter o ritmo de estudos para a disciplina.
    \end{swotlist}
  \end{tcolorbox}
  &
  % Oportunidades
  \begin{tcolorbox}[swotbox,
    title={OPORTUNIDADES (Opportunities)},
    colback=oporLight, colframe=opor, colbacktitle=opor,
    height=\rowheight]
    \begin{swotlist}
      \item \textbf{Biblioteca de excelência:} a do ISPTEC é uma das melhores do país, com amplo catálogo de recursos físicos e digitais.
      \item \textbf{Acesso a materiais complementares:} livros, artigos e outros materiais que ajudam a aprofundar os temas,incluindo livros escritos pelo docente, permitindo uma ferramenta extra para a compreensão da disciplina.
    \end{swotlist}
  \end{tcolorbox} \\[8pt]
  % Fraquezas
  \begin{tcolorbox}[swotbox,
    title={FRAQUEZAS (Weaknesses)},
    colback=fraqLight, colframe=fraq, colbacktitle=fraq,
    height=\rowheight]
    \begin{swotlist}
      \item \textbf{Complexidade da matéria:} conteúdo difícil de entender em alguns pontos, a quantidade de conhecimento de base necessário é maior comparada a outras disciplinas do mesmo semestre.
      \item \textbf{Risco de assimilação incompleta:} pode passar muita coisa sem que eu perceba que não entendi direito,e sem revisões constantes a situação agrava-se.
    \end{swotlist}
  \end{tcolorbox}
  &
  % Ameaças
  \begin{tcolorbox}[swotbox,
    title={AMEAÇAS (Threats)},
    colback=ameaLight, colframe=amea, colbacktitle=amea,
    height=\rowheight]
    \begin{swotlist}
      \item \textbf{Responsabilidades concorrentes:} trabalho, família e outras obrigações às vezes atrapalham os horários de estudo, falta de energia atrasa diversos projectos e sessões de estudo, imprevistos agravam a situação.
      \item \textbf{Gestão do tempo:} preciso organizar melhor, senão fica difícil concluir os trabalhos a tempo,priorizar ainda mais os meus estudos.
    \end{swotlist}
  \end{tcolorbox} \\
\end{tabularx}

\vfill

% ------------------ RODAPÉ ------------------
\noindent{\footnotesize\textcolor{axisColor!80}{%
  Análise centrada no sucesso individual na disciplina. %
  \hfill
  ISPTEC -- 2025%
}}

\end{document}